\section{Vyhodnocení klasifikace}\label{sec:ai-vyhodnoceni-klasifikace}

U~binární klasifikace porovnáváme skutečnou třídu \(y_i\in\set{0,1}\) s~předpovědí \(\hat y_i^{(t)}=h^{(t)}(\mathbf{x}_i)\). Mohou nastat čtyři výsledky:
\begin{itemize}
    \item správně pozitivní: \(y_i=1\) a~\(\hat y_i^{(t)}=1\);
    \item správně negativní: \(y_i=0\) a~\(\hat y_i^{(t)}=0\);
    \item falešně pozitivní: \(y_i=0\), ale \(\hat y_i^{(t)}=1\);
    \item falešně negativní: \(y_i=1\), ale \(\hat y_i^{(t)}=0\).
\end{itemize}

Jejich počty značíme
\[
    \TP_{\mathbf{X},\mathbf{Y}}^{(t)},
    \quad
    \TN_{\mathbf{X},\mathbf{Y}}^{(t)},
    \quad
    \FP_{\mathbf{X},\mathbf{Y}}^{(t)},
    \quad
    \FN_{\mathbf{X},\mathbf{Y}}^{(t)}.
\]
Například
\[
    \TP_{\mathbf{X},\mathbf{Y}}^{(t)}
    =
    \sizeof{
        \set{
            i:y_i=1,\ h^{(t)}(\mathbf{x}_i)=1
        }
    }.
\]
Ostatní tři počty definujeme analogicky. Je-li z~kontextu jasný soubor i~práh, budeme jejich dolní a~horní indexy někdy vynechávat.

\begin{figure}[ht]
    \centering
    \begin{tikzpicture}[every node/.style={font=\normalsize}]
    \node[task,fill=green!12,minimum width=43mm,minimum height=16mm]
        (tp) at (0,0) {\(\TP\)\\správně pozitivní};
    \node[task,fill=red!10,minimum width=43mm,minimum height=16mm]
        (fn) at (4.7,0) {\(\FN\)\\falešně negativní};
    \node[task,fill=red!10,minimum width=43mm,minimum height=16mm]
        (fp) at (0,-2) {\(\FP\)\\falešně pozitivní};
    \node[task,fill=green!12,minimum width=43mm,minimum height=16mm]
        (tn) at (4.7,-2) {\(\TN\)\\správně negativní};
    \node[above=7mm of tp] {předpověď \(1\)};
    \node[above=7mm of fn] {předpověď \(0\)};
    \node[left=5mm of tp,align=right] {skutečnost\\\(1\)};
    \node[left=5mm of fp,align=right] {skutečnost\\\(0\)};
\end{tikzpicture}

    \caption{Matice záměn pro binární klasifikaci.}
    \label{fig:ai-matice-zamen}
\end{figure}

\subsection{Správnost, přesnost, pokrytí a~\(F_1\)-míra}

\begin{definition}\label{def:ai-klasifikacni-metriky}
    Jsou-li příslušné jmenovatele nenulové, definujeme
    \begin{align*}
        \Accuracy_{\mathbf{X},\mathbf{Y}}^{(t)}
        &=
        \frac{\TP+\TN}{\TP+\TN+\FP+\FN},\\
        \Precision_{\mathbf{X},\mathbf{Y}}^{(t)}
        &=
        \frac{\TP}{\TP+\FP},\\
        \Recall_{\mathbf{X},\mathbf{Y}}^{(t)}
        &=
        \frac{\TP}{\TP+\FN},\\
        \Fscore_{1,\mathbf{X},\mathbf{Y}}^{(t)}
        &=
        \frac{2}{
            (\Precision_{\mathbf{X},\mathbf{Y}}^{(t)})^{-1}
            +
            (\Recall_{\mathbf{X},\mathbf{Y}}^{(t)})^{-1}
        }.
    \end{align*}
\end{definition}

\emph{Správnost} (anglicky \emph{accuracy}) odpovídá podílu všech správných odpovědí. \emph{Přesnost} (anglicky \emph{precision}) se ptá, kolika pozitivním předpovědím můžeme věřit. \emph{Pokrytí} (anglicky \emph{recall}) se ptá, jakou část skutečně pozitivních objektů jsme nalezli. \(F_1\)-míra je harmonický průměr přesnosti a~pokrytí a~je vysoká pouze tehdy, jsou-li vysoké obě hodnoty.

Dosazením definic dostaneme užitečný tvar
\[
    \Fscore_{1,\mathbf{X},\mathbf{Y}}^{(t)}
    =
    \frac{2\TP}{2\TP+\FP+\FN}.
\]

\begin{example}
    Pro
    \[
        \TP=36,\qquad
        \TN=44,\qquad
        \FP=8,\qquad
        \FN=12
    \]
    vychází
    \[
        \Accuracy=0{,}80,
        \quad
        \Precision\doteq0{,}82,
        \quad
        \Recall=0{,}75,
        \quad
        \Fscore_1\doteq0{,}78.
    \]
\end{example}

Správnost může být zavádějící u~nevyvážených tříd. Jestliže je mezi tisícem zpráv pouze deset podvodných, klasifikátor označující vždy běžnou zprávu dosáhne správnosti \(0{,}99\), ale jeho pokrytí podvodných zpráv je nulové.

\subsection{Vliv prahu a~neznámá data}

Snížení rozhodovacího prahu obvykle zvýší počet pozitivních předpovědí. Tím může růst pokrytí, ale zároveň přibývají falešně pozitivní výsledky a~klesá přesnost. Vhodný práh proto volíme podle důsledků obou typů chyb a~podle výsledků na validačních datech.

V~lékařském předběžném vyšetření může být závažnější přehlédnout nemocného člověka než poslat zdravého na další kontrolu; dává proto smysl upřednostnit vysoké pokrytí. U~automatického blokování účtů může být naopak falešně pozitivní rozhodnutí velmi nákladné, a~proto požadujeme vysokou přesnost. Jediná metrika bez kontextu tyto rozdíly nezachytí.

Změnou prahu získáme celou rodinu klasifikátorů ze stejného skórovacího modelu. Při jejich porovnávání sledujeme, jak se současně mění podíl správně pozitivních a~falešně pozitivních odpovědí. Tím lze hodnotit kvalitu samotného pořadí objektů podle skóre odděleně od volby jednoho konkrétního prahu.

\warningbox{Metriky uváděné po dokončení vývoje modelu počítáme na testovacích datech, která nebyla použita ani k~trénování, ani k~volbě prahu či jiných nastavení. Jinak získáme příliš optimistický odhad kvality.}
